% Este é o modelo do abnTeX2
% ---
% inserir lista de abreviaturas e siglas
% ---
%\begin{siglas}
%  \item[ABNT] Associação Brasileira de Normas Técnicas
%  \item[abnTeX] ABsurdas Normas para TeX
%\end{siglas}
% ---

%
% Outra opção é utilizar o pacote Acronym (mais funcional).
%

%\ac{SW}  	SW documents are beautifully typeset.
%\acf{SW} 	Scientific Word (SW) documents are beautifully typeset.
%\acs{SW} 	SW documents are beautifully typeset.
%\acl{SW} 	Scientific Word documents are beautifully typeset.

\pdfbookmark[0]{\listadesiglasname}{los}
\chapter*{\listadesiglasname}%

% Incluir no sumário
% \addcontentsline{toc}{chapter}{\listadesiglasname}%

\begin{acronym}
  \acro{ABNT}{Associação Brasileira de Normas Técnicas}
  \acro{API}{\textit{Application Programming Interface}}
  \acro{COSI}{Colegiado do Curso de Sistemas de Informação}
  \acro{CSV}{\textit{Comma-Separated Values}}
  \acro{DECSI}{Departamento de Computação e Sistemas}
  \acro{DOI}{\textit{Digital Object Identifier}}
  \acro{GEXF}{\textit{Graph Exchange XML Format}}
  \acro{HTML}{\textit{HyperText Markup Language}}
  \acro{ICEA}{Instituto de Ciências Exatas e Aplicadas}
  \acro{IMC}{Índice de Massa Corporal}
  \acro{JM}{João Monlevade}
  \acro{JSON}{\textit{JavaScript Object Notation}}
  \acro{NBA}{\textit{National Basketball Association}}
  \acro{NOC}{\textit{National Olympic Committee}}
  \acro{SI}{Sistemas de Informação}
  \acro{TCC}{Trabalho de Conclusão de Curso}
  \acro{UFOP}{Universidade Federal de Ouro Preto}
  \acro{XML}{\textit{eXtensible Markup Language}}
\end{acronym}
